% ---------------------------------------------------------------------------
% Motivação FOTA: o instante de início da transferência muda tudo.
% Ilustração original, inspirada na Fig. 1 de Andrade et al. (2019).
%
% Cada barra é a FRAÇÃO DA ATUALIZAÇÃO baixada naquele dia; a altura útil da
% faixa de um cenário vale 100%. Assim os dois cenários somam o mesmo total e
% as quatro parcelas são diretamente comparáveis. Para reajustar o contraste,
% basta mexer nas quatro frações abaixo.
%
% Os `pics` do veículo são desenhados em cm (e não em unidades da figura) para
% não se deformarem quando a razão y/x muda.
% ---------------------------------------------------------------------------
\begin{tikzpicture}[
  x=1.25cm, y=1.78cm,
  font=\footnotesize,
  cell/.style={draw=black!35, line width=0.4pt},
  data/.style={fill=mydarkblue1, draw=mydarkblue1!55!black, line width=0.4pt},
  daychip/.style={draw, rounded corners=1.5pt, line width=0.5pt,
                  minimum width=0.82cm, minimum height=0.32cm,
                  inner sep=1pt, font=\footnotesize},
  rowlabel/.style={anchor=east, align=right, font=\footnotesize},
  verdict/.style={anchor=west, align=left, font=\footnotesize\bfseries},
]

% --- parâmetros -------------------------------------------------------------
\def\rowh{1.00}        % altura útil da faixa de um cenário = 100% da atualização
\def\barw{0.76}        % largura da barra, em frações da largura de um dia
\def\ya{1.95}          % base da faixa do cenário (a)
\def\yb{0.62}          % base da faixa do cenário (b)

\def\fracAseg{0.50}    % (a) metade na segunda ...
\def\fracAqua{0.50}    %     ... e metade na quarta -- ambas no pico
\def\fracBsex{0.80}    % (b) quase tudo na sexta, fora do pico ...
\def\fracBseg{0.20}    %     ... e só o resto na segunda

% --- barra de dados: #1 dia, #2 base da faixa, #3 fração da atualização -----
\newcommand{\dlbar}[3]{%
  \fill[data] (#1+0.5-\barw/2, #2) rectangle ++(\barw, #3*\rowh);
}

% --- faixa de um cenário: #1 base ------------------------------------------
\newcommand{\dayrow}[1]{%
  \draw[cell] (0,#1) rectangle (8,#1+\rowh);
  \foreach \i in {1,...,7} { \draw[cell] (\i,#1) -- (\i,#1+\rowh); }
}

% --- pics: veículo conectado (carro + arcos de sinal) -----------------------
\tikzset{
  carro/.pic={
    \fill[mydarkblue1!85] (-0.33cm,0.02cm) -- (-0.33cm,0.17cm)
      -- (-0.15cm,0.17cm) -- (-0.04cm,0.32cm) -- (0.18cm,0.32cm)
      -- (0.21cm,0.17cm) -- (0.33cm,0.17cm) -- (0.33cm,0.02cm) -- cycle;
    \fill[mydarkblue1!85] (-0.18cm,0cm) circle (0.07cm);
    \fill[mydarkblue1!85] (0.18cm,0cm) circle (0.07cm);
  },
  sinal/.pic={
    \draw[mydarkblue1!85, line width=0.5pt] (-0.12cm,0cm) arc (180:0:0.12cm);
    \draw[mydarkblue1!85, line width=0.5pt] (-0.23cm,0cm) arc (180:0:0.23cm);
    \fill[mydarkblue1!85] (0cm,-0.04cm) circle (0.035cm);
  }
}

% ===========================================================================
% Faixa de pico da rede (segunda a quinta), atrás de tudo
% ===========================================================================
\fill[mylightorange1, opacity=0.45] (3,\yb-0.14) rectangle (7,3.44);

% Colchete dos períodos de pico, no estilo da Fig. 1 do artigo
\draw[mydarkorange1, line width=0.8pt]
  (3.10,3.98) -- (6.90,3.98)
  (3.10,3.98) -- (3.10,3.90)
  (6.90,3.98) -- (6.90,3.90);
\node[mydarkorange1, font=\footnotesize\bfseries, anchor=south] at (5,4.02)
  {per\'iodos de pico da rede};

% --- dias -------------------------------------------------------------------
\foreach \i/\d in {0/Sex, 1/S\'ab, 2/Dom, 7/Sex} {
  \node[daychip, fill=mygray!35, draw=black!45] at (\i+0.5, 3.68) {\d};
}
\foreach \i/\d in {3/Seg, 4/Ter, 5/Qua, 6/Qui} {
  \node[daychip, fill=mylightorange1, draw=mydarkorange1,
        text=mydarkorange1!70!black] at (\i+0.5, 3.68) {\d};
}

% --- disponibilidade do veículo (a mesma nos dois cenários) ----------------
\node[rowlabel, text=mydarkblue1!80!black] at (-0.16, 3.19) {ve\'iculo na rede};
\foreach \i in {0,3,5,7} {
  \path (\i+0.5,3.10) pic {carro};
  \path (\i+0.5,3.30) pic {sinal};
}

% ===========================================================================
% Cenário (a): início na segunda-feira -- tudo cai no pico
% ===========================================================================
\node[rowlabel] at (-0.16, \ya+\rowh/2) {\textbf{(a)} in\'icio\\na segunda};
\dayrow{\ya}
\dlbar{3}{\ya}{\fracAseg}
\dlbar{5}{\ya}{\fracAqua}
\node[verdict, mydarkorange1] at (8.16, \ya+\rowh/2)
  {toda a\\transfer\^encia\\cai no pico};

% ===========================================================================
% Cenário (b): início na sexta-feira -- quase tudo fora do pico
% ===========================================================================
\node[rowlabel] at (-0.16, \yb+\rowh/2) {\textbf{(b)} in\'icio\\na sexta};
\dayrow{\yb}
\dlbar{0}{\yb}{\fracBsex}
\dlbar{3}{\yb}{\fracBseg}
\node[verdict, mygreen!55!black] at (8.16, \yb+\rowh/2)
  {quase tudo\\fora do pico};

% --- legenda ---------------------------------------------------------------
\fill[data] (0.06,0.06) rectangle (0.50,0.24);
\node[anchor=west, font=\footnotesize] at (0.60,0.15)
  {dados transferidos (fra\c{c}\~ao da atualiza\c{c}\~ao)};

\end{tikzpicture}
